\begin{recipe}
[% 
    preparationtime = {\unit[30]{min}},
    bakingtime={2 x \unit[8--12]{min}},
    bakingtemperature={\protect\bakingtemperature{
        topbottomheat=\unit[250]{\textcelcius}}},
    portion = {\portion{2}},
    source = {Adrian Hieber},
    rating = {3},
    calory = {915 kcal | 44g Protein | 95g KH | 40g Fett},
]
{Flammkuchen}
    
    % \graph
    % {% pictures
    %     small=pic/flammkuchen_small,
    %     big=pic/flammkuchen_big
    % }
    
    \introduction{%
        Selbstgemachter Flammkuchen geht einfacher als man denkt. 
        Dünner Teig, cremiger Belag und richtig heiß gebacken -- 
        perfekt für einen entspannten Abend oder wenn Besuch da ist.
        
    }
    
    \ingredients{%
        215 g & Mehl\\
        1{,}5 g & Trockenhefe\\
        5 g & Salz\\
        1{,}5 g & Zucker\\
        10 g & Olivenöl\\
        130 g & Wasser\\
        125 g & Creme (z.B. vegane Creme oder Creme fraîche)\\
        50 g & Joghurt oder Sahne\\
        \nicefrac{1}{2} & Rote Zwiebel\\
        175 g & Räuchertofu\\
        1 Msp. & Muskat\\
        1 EL & Schnittlauch
    }
    
    \preparation{%
        \begin{enumerate}
            \item \textbf{Teig:} Mehl, Hefe, Salz und Zucker vermischen. 
                  Wasser und Olivenöl hinzufügen und zu einem geschmeidigen Teig kneten.
            
            \item Teig halbieren und auf Backpapier sehr dünn ausrollen.
            
            \item \textbf{Belag:} Zwiebel in feine Streifen schneiden. 
                  Räuchertofu in kleine Würfel schneiden.
            
            \item Creme, Joghurt/Sahne, Schnittlauch und Muskat verrühren.
            
            \item Flammkuchen dünn mit der Creme bestreichen und mit Tofu und Zwiebeln belegen.
            
            \item Bei \unit[250]{\textcelcius} etwa 8--12 Minuten backen, 
                  bis der Rand goldbraun ist. Vorgang mit dem zweiten Flammkuchen wiederholen.
        \end{enumerate}
    }
    
    \hint{%
        Je dünner der Teig, desto besser wird der Flammkuchen. 
        Der Ofen sollte wirklich gut vorgeheizt sein.
    }
    
\end{recipe}